% Discussion — Paper 2, Medical Physics submission
% Humanised v1, 2026-05-25. Formal-academic-active register, citation-bound.
% Navigates: (1) honest H1 rejection, (2) H1' reframe with exploratory disclosure,
% (3) size-confound finding as primary methodological contribution,
% (4) explicit Pai 2024 (peer-reviewed) differentiation,
% (5) Limitations including the single-FM caveat.

\subsection{Summary of findings}

We compared the propagation of segmentation-mask perturbation through pretrained foundation-model embeddings (BiomedCLIP)~\cite{zhang2024biomedclip} and IBSI-aligned handcrafted radiomic features~\cite{vangriethuysen2017,zwanenburg2020} extracted in matched 2D and 3D configurations from the same lung-CT nodule cohort (LIDC-IDRI~\cite{armato2011}; 310 balanced nodules with up to four-reader natural segmentation variability per nodule, augmented by seven controlled morphological perturbations and one Dice-calibrated synthetic mask per nodule). The four pre-specified hypotheses produced a mixed picture. The hypothesis that foundation-model embeddings exhibit a higher proportion of stable (ICC $\geq 0.75$) features than handcrafted radiomics was rejected; both paradigms achieved approximately 90\% stable feature rates. The empirical distributions of ICC values, however, differed substantially between paradigms: radiomics produced a bimodal distribution with a large excellent-stratum mode and a small unstable tail, whereas BiomedCLIP produced a unimodal distribution concentrated in the moderate-to-good stratum with no poor (ICC $<$ 0.50) dimensions. The downstream malignancy-classification analysis showed that ICC-based stability filtering had essentially no effect on radiomics AUC, while it improved BiomedCLIP AUC modestly and improved its calibration substantially. The pre-specified diameter-stratified subgroup analysis revealed that the apparent aggregate AUC superiority of handcrafted radiomics over BiomedCLIP (0.95 versus 0.90) was substantially attributable to a size-malignancy confound endemic to LIDC, rather than to a paradigm-level feature-quality difference: within the matched-size 6--10~mm stratum, the three paradigms achieved statistically indistinguishable AUC values around 0.77.

\subsection{Comparison to the peer-reviewed literature}

The peer-reviewed precedent for foundation-model stability analysis in cancer imaging is Pai~\textit{et al.}~2024~\cite{pai2024}, who introduced FMCIB --- a three-dimensional contrastively pretrained ResNet50 --- and reported ICC = 0.984 for its feature-based predictions on the RIDER same-day repeat-scan dataset, with significantly higher stability than a supervised CNN baseline under three-dimensional Gaussian seed-point jitter ($\sigma^2 = 16$ voxels per dimension) on the LUNG1 cohort. Pai~\textit{et al.}~2024 established that domain-pretrained 3D foundation models can be more stable than supervised CNN baselines under scanner-repeat and seed-point variability. The present study extends this peer-reviewed evidence base in three concrete directions. First, the seed-point-jitter perturbation is replaced with controlled segmentation-mask perturbation, which is the dominant noise source in handcrafted radiomic pipelines and the variability mode that radiomics practitioners actually filter for via ICC~\cite{pavic2018,vantimmeren2020}. Second, an IBSI-aligned PyRadiomics comparator is extracted from the same lesions under the same perturbations, enabling a true head-to-head paradigm comparison rather than an FM-internal benchmark. Third, ICC is computed per dimension rather than as a single-output statistic across the embedding, which reveals the distribution-shape divergence (\S\,3.3) that single-output stability statistics necessarily conceal. A contemporaneous preprint by the same group extends the FMCIB benchmark to ten 3D foundation models using global cosine similarity; this work has not undergone peer review and is referenced here only for completeness, without being treated as established precedent. Cosma~\textit{et al.}~2025~\cite{cosma2025}, working independently on breast MRI for triple-negative breast cancer subtype prediction, reported the counterintuitive finding that ICC-based filtering of radiomic features can exclude features with predictive value; the H3 result reported here extends and corroborates that finding for both handcrafted radiomics and foundation-model embeddings on lung CT.

\subsection{Mechanistic interpretation of the distribution-shape divergence}

The most striking finding of this study is the substantial difference in ICC distribution shape between paradigms (Kolmogorov--Smirnov $D = 0.72$, $p < 10^{-43}$ between BiomedCLIP and PyRadiomics-2D). Two mechanistic explanations are consistent with this pattern. First, the SimCLR-style contrastive pretraining objective underlying BiomedCLIP~\cite{zhang2024biomedclip} enforces approximately uniform invariance to image augmentations across every embedding dimension: the loss function rewards the encoder for producing nearly identical vectors for two augmented views of the same image, and this pressure operates uniformly across all 512 dimensions of the output. The empirical consequence is a unimodal distribution concentrated in the moderate-to-good stability band with no catastrophically unstable dimensions. Second, IBSI radiomic features have heterogeneous mathematical sensitivities to mask perturbation by construction~\cite{zwanenburg2020}: features such as \texttt{firstorder\_Mean}, \texttt{firstorder\_TotalEnergy}, \texttt{shape\_VoxelVolume}, and \texttt{shape\_MeshVolume} are mathematically near-invariant to small boundary perturbations, whereas certain GLSZM, GLDM, and high-frequency wavelet-derived texture features are inherently sensitive to mask edges and produce the small unstable tail in the radiomics distribution. The bimodal radiomics distribution thus reflects the underlying bimodality of the feature definitions: a robust core of intensity-and-shape descriptors plus a noise-sensitive periphery of high-order texture descriptors. The two paradigms differ not primarily in \emph{how much} stability they provide, but in \emph{which dimensions} carry the stability --- a methodologically important distinction that aggregate proportion-stable metrics conceal.

A second mechanistic observation arises from the calibration analysis (\S\,3.4). BiomedCLIP's raw 512-dimensional embeddings produced poorly calibrated probability outputs (Brier score 0.144, expected calibration error 0.129) even at strong discrimination (AUC 0.898). ICC-based filtering to the 84 dimensions with ICC $\geq 0.85$ reduced both Brier and ECE by $\geq 20$\% while simultaneously improving discrimination. This is consistent with an interpretation in which the moderately-stable BiomedCLIP dimensions encode perturbation-driven noise that adds calibration uncertainty without contributing malignancy signal: when these dimensions are removed, the L2-regularised logistic regression no longer needs to balance contributions from unstable noise dimensions, and the resulting probability outputs are sharper and better calibrated. To our knowledge this calibration-improvement effect under per-dimension ICC selection has not been previously reported for foundation-model embeddings, and it is clinically actionable for any deployment pipeline that reports probability outputs rather than ranked predictions.

\subsection{Size-confound finding and implications for the FM-radiomics benchmarking literature}

The size-stratified subgroup analysis (\S\,3.5; Figure~4) is, in our assessment, the most methodologically consequential finding of this study. The unfiltered aggregate AUC ranking --- handcrafted radiomics at 0.95 versus BiomedCLIP at 0.90 --- was substantially driven by the structural correlation between nodule diameter and malignancy in LIDC-IDRI, rather than by paradigm-level feature-quality differences. Handcrafted radiomic features include explicit shape descriptors (\texttt{shape\_VoxelVolume}, \texttt{shape\_MaximumDiameter}, \texttt{shape\_MeshVolume}) that encode lesion size directly~\cite{vangriethuysen2017,zwanenburg2020}. The BiomedCLIP preprocessing pipeline used here --- bounding-box cropping with 50\% padding followed by resampling to $224 \times 224$ pixels --- normalises absolute size out of the input by design, as is standard for vision-transformer inference. The aggregate AUC comparison therefore confounded ``radiomics has explicit size features'' with ``radiomics produces higher-quality features for malignancy classification.'' Within the matched-size 6--10~mm stratum, the AUC gap dissolved to within the 95\% confidence-interval overlap (0.762 / 0.774 / 0.758 for the three paradigms); within the matched-size 10--20~mm stratum a residual $\sim 0.095$ AUC advantage for PyRadiomics-3D persisted with overlapping confidence intervals, plausibly reflecting genuine sub-stratum information in three-dimensional shape descriptors that the normalised single-slice BiomedCLIP input cannot recover.

The methodological implication for the wider FM-versus-radiomics benchmarking literature is direct: any benchmark study performed on a cohort with structural size-class imbalance will systematically inflate handcrafted radiomics' apparent advantage relative to foundation models with size-normalising preprocessing. Future FM-radiomics comparisons should report size-stratified AUC values as the primary discrimination metric and treat aggregate AUC as a confounded secondary statistic. The apparent heterogeneity of published FM-versus-radiomics comparisons across cohorts may, in substantial part, be a consequence of unmeasured size-class composition differences across those cohorts.

\subsection{Limitations}

Several limitations of this work should be acknowledged. First, the foundation-model arm comprises a single model (BiomedCLIP)~\cite{zhang2024biomedclip}; RadImageNet-ResNet50~\cite{mei2022radimagenet} was planned as a second paradigm but the required weight access was not granted within the study timeframe. The findings reported here are therefore strictly BiomedCLIP-specific, and generalisation to other foundation-model classes --- including 3D models such as FMCIB~\cite{pai2024} and segmentation-pretrained models such as MedSAM~\cite{ma2024medsam} --- requires the planned multi-FM amendment. Second, the entire study was performed on a single anatomical site (lung), a single modality (computed tomography), and a single cohort (LIDC-IDRI); cross-site, cross-modality, and cross-cohort generalisation are deferred to a planned external-validation study. Third, the 6~mm diameter floor reflects a pre-registered fallback cascade from the originally pre-registered 10~mm floor~\cite{macmahon2017}; the sensitivity analysis at 10~mm preserves protocol-fidelity reporting but reduces $N$ to 129 with substantial loss of statistical power. Fourth, the $>20$~mm stratum was not evaluable due to LIDC's structural scarcity of large benign nodules, a corpus-design limitation rather than an analytical one. Fifth, all feature extraction was performed on CPU; the BiomedCLIP two-dimensional inference protocol was chosen partly for this constraint, and a three-dimensional FM extension on GPU-enabled compute is a natural follow-on. Sixth, the L2-regularised logistic regression chosen as the downstream classifier~\cite{hosmer2013} was selected for its bit-exact reproducibility properties and its standard use in the radiomics ICC-filter literature; alternative classifiers (gradient-boosting, random forests, regularised survival models) may exhibit different sensitivities to feature-set composition.

\subsection{Future work}

The most immediate planned extension is the multi-FM amendment incorporating RadImageNet-ResNet50~\cite{mei2022radimagenet} as a second foundation-model paradigm once weight access is granted, with the MedSAM image-encoder~\cite{ma2024medsam} as a robust fallback. A natural three-dimensional FM extension using FMCIB~\cite{pai2024} and additional 3D encoders on GPU-enabled compute would address the 2D--3D dimensionality limitation. A multi-cohort external-validation extension to NSCLC-Radiomics, NSCLC-Radiogenomics, and an Australian rural-cohort dataset would address the cross-site generalisation limitation. Finally, the calibration-improvement effect observed for BiomedCLIP (\S\,3.4) warrants a dedicated study replicating the per-dimension ICC characterisation in additional foundation models, to test whether the effect is a general property of contrastively pretrained embeddings or specific to BiomedCLIP's particular training corpus and architecture.
